
\documentclass{article}
\usepackage{colortbl}
\usepackage{makecell}
\usepackage{multirow}
\usepackage{supertabular}

\begin{document}

\newcounter{utterance}

\centering \large Interaction Transcript for game `matchit\_ascii', experiment `similar\_grid\_2', episode 4 with deepseek{-}v4{-}pro.
\vspace{24pt}

{ \footnotesize  \setcounter{utterance}{1}
\setlength{\tabcolsep}{0pt}
\begin{supertabular}{c@{$\;$}|p{.15\linewidth}@{}p{.15\linewidth}p{.15\linewidth}p{.15\linewidth}p{.15\linewidth}p{.15\linewidth}}
    \# & $\;$A & \multicolumn{4}{c}{Game Master} & $\;\:$B\\
    \hline

    \theutterance \stepcounter{utterance}  
    & & \multicolumn{4}{p{0.6\linewidth}}{
        \cellcolor[rgb]{0.9,0.9,0.9}{
            \makecell[{{p{\linewidth}}}]{
                \texttt{\tiny{[P1$\langle$GM]}}
                \texttt{Sudjeluješ u igri pogađanja u suradnji.} \\
\texttt{Cilj je otkriti jesu li ova mreža i druga slika koju samo ja mogu vidjeti iste.} \\
\texttt{Molim te, prvo opiši svoju mrežu. Zatim ću ja dati svoj opis i možemo si međusobno postavljati pitanja o mrežama kako bismo utvrdili jesu li iste. Sada započni svoj kratki opis mreže s "OPIS:" nakon čega slijedi opis. Ne dodaj ništa drugo. Ovo je tvoja mreža:} \\
\texttt{X X X X X} \\
\texttt{X ▢ ▢ ▢ X} \\
\texttt{X ▢ ▢ ▢ X} \\
\texttt{X ▢ ▢ ▢ X} \\
\texttt{X X X X X} \\
            }
        }
    }
    & \\ \\

    \theutterance \stepcounter{utterance}  
    & \multicolumn{4}{p{0.6\linewidth}}{
        \cellcolor[rgb]{0.9,0.9,0.9}{
            \makecell[{{p{\linewidth}}}]{
                \texttt{\tiny{[P1$\rangle$GM]}}
                \texttt{OPIS:} \\
\texttt{Mreža je veličine 5x5.} \\
\texttt{Vanjski rub se sastoji isključivo od znakova X, tvoreći okvir.} \\
\texttt{Unutrašnjost mreže (3x3 polja) sastoji se isključivo od praznih kvadrata (▢).} \\
            }
        }
    }
    & & \\ \\

    \theutterance \stepcounter{utterance}  
    & & & \multicolumn{2}{p{0.3\linewidth}}{
        \cellcolor[rgb]{0.9,0.9,0.9}{
            \makecell[{{p{\linewidth}}}]{
                \texttt{\tiny{[GM$|$GM]}}
                \texttt{continue} \\
            }
        }
    }
    & & \\ \\

    \theutterance \stepcounter{utterance}  
    & & \multicolumn{4}{p{0.6\linewidth}}{
        \cellcolor[rgb]{0.9,0.9,0.9}{
            \makecell[{{p{\linewidth}}}]{
                \texttt{\tiny{[GM$\rangle$P2]}}
                \texttt{Sudjeluješ u igri pogađanja u suradnji.} \\
\texttt{Cilj je otkriti jesu li ova mreža i druga slika koju samo ja mogu vidjeti iste.} \\
\texttt{Molim te, prvo opiši svoju mrežu. Zatim ću ja dati svoj opis i možemo si međusobno postavljati pitanja o mrežama kako bismo utvrdili jesu li iste. Sada započni svoj kratki opis mreže s "OPIS:" nakon čega slijedi opis. Ne dodaj ništa drugo. Ovo je tvoja mreža:} \\
\texttt{X ▢ X X X} \\
\texttt{X ▢ ▢ ▢ X} \\
\texttt{X ▢ ▢ ▢ X} \\
\texttt{X ▢ ▢ ▢ X} \\
\texttt{X ▢ X X X} \\
            }
        }
    }
    & \\ \\

    \theutterance \stepcounter{utterance}  
    & & & \multicolumn{4}{p{0.6\linewidth}}{
        \cellcolor[rgb]{0.9,0.9,0.9}{
            \makecell[{{p{\linewidth}}}]{
                \texttt{\tiny{[GM$\langle$P2]}}
                \texttt{OPIS: Mreža je dimenzija 5x5. Rubovi su ispunjeni znakovima X, osim četiri prazna polja (▢) na sredini svakog ruba. Unutrašnjost (3x3) je ispunjena praznim poljima (▢), osim središnjeg retka i stupca? Čekaj, da preciznije opišem: } \\
\texttt{Gornji red: X, prazno, X, X, X} \\
\texttt{Drugi red: X, prazno, prazno, prazno, X} \\
\texttt{Treći red: X, prazno, prazno, prazno, X} \\
\texttt{Četvrti red: X, prazno, prazno, prazno, X} \\
\texttt{Donji red: X, prazno, X, X, X} \\
\\ 
\texttt{Uglavnom, to je okvir od X{-}ova s jednim praznim poljem na sredini gornjeg i donjeg ruba, a unutrašnjost je prazna osim što su lijevi i desni rubovi potpuno X{-}ovi? Oprosti, ponovno gledam:} \\
\texttt{Red 1: X, prazno, X, X, X} \\
\texttt{Red 2: X, prazno, prazno, prazno, X} \\
\texttt{Red 3: X, prazno, prazno, prazno, X} \\
\texttt{Red 4: X, prazno, prazno, prazno, X} \\
\texttt{Red 5: X, prazno, X, X, X} \\
\\ 
\texttt{Znači, oblik podsjeća na slovo "C" okrenuto udesno ili možda utor? Samo jedno prazno polje na gornjem rubu (stupac 2) i jedno na donjem rubu (stupac 2). Lijevi i desni rub su potpuno X{-}ovi. Unutrašnjost (redovi 2{-}4, stupci 2{-}4) su sva prazna polja.} \\
            }
        }
    }
     \\ \\

    \theutterance \stepcounter{utterance}  
    & & & \multicolumn{2}{p{0.3\linewidth}}{
        \cellcolor[rgb]{0.9,0.9,0.9}{
            \makecell[{{p{\linewidth}}}]{
                \texttt{\tiny{[GM$|$GM]}}
                \texttt{continue} \\
            }
        }
    }
    & & \\ \\

    \theutterance \stepcounter{utterance}  
    & & \multicolumn{4}{p{0.6\linewidth}}{
        \cellcolor[rgb]{0.9,0.9,0.9}{
            \makecell[{{p{\linewidth}}}]{
                \texttt{\tiny{[GM$\rangle$P2]}}
                \texttt{Ovo je mojOPIS:} \\
\texttt{Mreža je veličine 5x5.} \\
\texttt{Vanjski rub se sastoji isključivo od znakova X, tvoreći okvir.} \\
\texttt{Unutrašnjost mreže (3x3 polja) sastoji se isključivo od praznih kvadrata (▢).} \\
\texttt{Sada postavi pitanje kako bi saznao nove aspekte moje mreže koji se možda razlikuju od tvoje mreže. Započni s "PITANJE:" i ne dodavaj ništa drugo.} \\
            }
        }
    }
    & \\ \\

    \theutterance \stepcounter{utterance}  
    & & & \multicolumn{4}{p{0.6\linewidth}}{
        \cellcolor[rgb]{0.9,0.9,0.9}{
            \makecell[{{p{\linewidth}}}]{
                \texttt{\tiny{[GM$\langle$P2]}}
                \texttt{PITANJE: Nalazi li se u tvom gornjem redu (1. red) prazan kvadrat (▢) na poziciji stupca 2?} \\
            }
        }
    }
     \\ \\

    \theutterance \stepcounter{utterance}  
    & & & \multicolumn{2}{p{0.3\linewidth}}{
        \cellcolor[rgb]{0.9,0.9,0.9}{
            \makecell[{{p{\linewidth}}}]{
                \texttt{\tiny{[GM$|$GM]}}
                \texttt{continue} \\
            }
        }
    }
    & & \\ \\

    \theutterance \stepcounter{utterance}  
    & & \multicolumn{4}{p{0.6\linewidth}}{
        \cellcolor[rgb]{0.9,0.9,0.9}{
            \makecell[{{p{\linewidth}}}]{
                \texttt{\tiny{[P1$\langle$GM]}}
                \texttt{Ovo je mojOPIS: Mreža je dimenzija 5x5. Rubovi su ispunjeni znakovima X, osim četiri prazna polja (▢) na sredini svakog ruba. Unutrašnjost (3x3) je ispunjena praznim poljima (▢), osim središnjeg retka i stupca? Čekaj, da preciznije opišem: } \\
\texttt{Gornji red: X, prazno, X, X, X} \\
\texttt{Drugi red: X, prazno, prazno, prazno, X} \\
\texttt{Treći red: X, prazno, prazno, prazno, X} \\
\texttt{Četvrti red: X, prazno, prazno, prazno, X} \\
\texttt{Donji red: X, prazno, X, X, X} \\
\\ 
\texttt{Uglavnom, to je okvir od X{-}ova s jednim praznim poljem na sredini gornjeg i donjeg ruba, a unutrašnjost je prazna osim što su lijevi i desni rubovi potpuno X{-}ovi? Oprosti, ponovno gledam:} \\
\texttt{Red 1: X, prazno, X, X, X} \\
\texttt{Red 2: X, prazno, prazno, prazno, X} \\
\texttt{Red 3: X, prazno, prazno, prazno, X} \\
\texttt{Red 4: X, prazno, prazno, prazno, X} \\
\texttt{Red 5: X, prazno, X, X, X} \\
\\ 
\texttt{Znači, oblik podsjeća na slovo "C" okrenuto udesno ili možda utor? Samo jedno prazno polje na gornjem rubu (stupac 2) i jedno na donjem rubu (stupac 2). Lijevi i desni rub su potpuno X{-}ovi. Unutrašnjost (redovi 2{-}4, stupci 2{-}4) su sva prazna polja.} \\
\texttt{PITANJE: Nalazi li se u tvom gornjem redu (1. red) prazan kvadrat (▢) na poziciji stupca 2?Započni svoj odgovor s "ODGOVOR:" i nemoj dodati ništa drugo.} \\
            }
        }
    }
    & \\ \\

    \theutterance \stepcounter{utterance}  
    & \multicolumn{4}{p{0.6\linewidth}}{
        \cellcolor[rgb]{0.9,0.9,0.9}{
            \makecell[{{p{\linewidth}}}]{
                \texttt{\tiny{[P1$\rangle$GM]}}
                \texttt{ODGOVOR: Ne, u mom gornjem redu nema praznih kvadrata; svi su X.} \\
            }
        }
    }
    & & \\ \\

    \theutterance \stepcounter{utterance}  
    & & & \multicolumn{2}{p{0.3\linewidth}}{
        \cellcolor[rgb]{0.9,0.9,0.9}{
            \makecell[{{p{\linewidth}}}]{
                \texttt{\tiny{[GM$|$GM]}}
                \texttt{continue} \\
            }
        }
    }
    & & \\ \\

    \theutterance \stepcounter{utterance}  
    & & \multicolumn{4}{p{0.6\linewidth}}{
        \cellcolor[rgb]{0.9,0.9,0.9}{
            \makecell[{{p{\linewidth}}}]{
                \texttt{\tiny{[P1$\langle$GM]}}
                \texttt{Sada postavi pitanje kako bi saznao nove aspekte moje mreže koji se možda razlikuju od tvoje mreže. Započni s "PITANJE:" i ne dodavaj ništa drugo.} \\
            }
        }
    }
    & \\ \\

    \theutterance \stepcounter{utterance}  
    & \multicolumn{4}{p{0.6\linewidth}}{
        \cellcolor[rgb]{0.9,0.9,0.9}{
            \makecell[{{p{\linewidth}}}]{
                \texttt{\tiny{[P1$\rangle$GM]}}
                \texttt{PITANJE: Nalaze li se u tvom lijevom rubnom stupcu (stupac 1) isključivo znakovi X u svih pet redova?} \\
            }
        }
    }
    & & \\ \\

    \theutterance \stepcounter{utterance}  
    & & & \multicolumn{2}{p{0.3\linewidth}}{
        \cellcolor[rgb]{0.9,0.9,0.9}{
            \makecell[{{p{\linewidth}}}]{
                \texttt{\tiny{[GM$|$GM]}}
                \texttt{continue} \\
            }
        }
    }
    & & \\ \\

    \theutterance \stepcounter{utterance}  
    & & \multicolumn{4}{p{0.6\linewidth}}{
        \cellcolor[rgb]{0.9,0.9,0.9}{
            \makecell[{{p{\linewidth}}}]{
                \texttt{\tiny{[GM$\rangle$P2]}}
                \texttt{ODGOVOR: Ne, u mom gornjem redu nema praznih kvadrata; svi su X.} \\
\texttt{PITANJE: Nalaze li se u tvom lijevom rubnom stupcu (stupac 1) isključivo znakovi X u svih pet redova?Započni svoj odgovor s "ODGOVOR:" i nemoj dodati ništa drugo.} \\
            }
        }
    }
    & \\ \\

    \theutterance \stepcounter{utterance}  
    & & & \multicolumn{4}{p{0.6\linewidth}}{
        \cellcolor[rgb]{0.9,0.9,0.9}{
            \makecell[{{p{\linewidth}}}]{
                \texttt{\tiny{[GM$\langle$P2]}}
                \texttt{ODGOVOR: Da, u mom lijevom rubnom stupcu (stupac 1) nalaze se isključivo znakovi X u svih pet redova.} \\
\texttt{PITANJE: Ima li tvoj donji red (red 5) prazan kvadrat (▢) na poziciji stupca 2?} \\
            }
        }
    }
     \\ \\

    \theutterance \stepcounter{utterance}  
    & & & \multicolumn{2}{p{0.3\linewidth}}{
        \cellcolor[rgb]{0.9,0.9,0.9}{
            \makecell[{{p{\linewidth}}}]{
                \texttt{\tiny{[GM$|$GM]}}
                \texttt{continue} \\
            }
        }
    }
    & & \\ \\

    \theutterance \stepcounter{utterance}  
    & & \multicolumn{4}{p{0.6\linewidth}}{
        \cellcolor[rgb]{0.9,0.9,0.9}{
            \makecell[{{p{\linewidth}}}]{
                \texttt{\tiny{[GM$\rangle$P2]}}
                \texttt{Sada postavi pitanje kako bi saznao nove aspekte moje mreže koji se možda razlikuju od tvoje mreže. Započni s "PITANJE:" i ne dodavaj ništa drugo.} \\
            }
        }
    }
    & \\ \\

    \theutterance \stepcounter{utterance}  
    & & & \multicolumn{4}{p{0.6\linewidth}}{
        \cellcolor[rgb]{0.9,0.9,0.9}{
            \makecell[{{p{\linewidth}}}]{
                \texttt{\tiny{[GM$\langle$P2]}}
                \texttt{PITANJE: Ima li tvoj desni rubni stupac (stupac 5) isključivo znakove X u svih pet redova?} \\
            }
        }
    }
     \\ \\

    \theutterance \stepcounter{utterance}  
    & & & \multicolumn{2}{p{0.3\linewidth}}{
        \cellcolor[rgb]{0.9,0.9,0.9}{
            \makecell[{{p{\linewidth}}}]{
                \texttt{\tiny{[GM$|$GM]}}
                \texttt{continue} \\
            }
        }
    }
    & & \\ \\

    \theutterance \stepcounter{utterance}  
    & & \multicolumn{4}{p{0.6\linewidth}}{
        \cellcolor[rgb]{0.9,0.9,0.9}{
            \makecell[{{p{\linewidth}}}]{
                \texttt{\tiny{[P1$\langle$GM]}}
                \texttt{ODGOVOR: Da, u mom lijevom rubnom stupcu (stupac 1) nalaze se isključivo znakovi X u svih pet redova.} \\
\texttt{PITANJE: Ima li tvoj donji red (red 5) prazan kvadrat (▢) na poziciji stupca 2?} \\
\texttt{PITANJE: Ima li tvoj desni rubni stupac (stupac 5) isključivo znakove X u svih pet redova?Započni svoj odgovor s "ODGOVOR:" i nemoj dodati ništa drugo.} \\
            }
        }
    }
    & \\ \\

    \theutterance \stepcounter{utterance}  
    & \multicolumn{4}{p{0.6\linewidth}}{
        \cellcolor[rgb]{0.9,0.9,0.9}{
            \makecell[{{p{\linewidth}}}]{
                \texttt{\tiny{[P1$\rangle$GM]}}
                \texttt{ODGOVOR: } \\
\texttt{Ne, moj donji red nema prazan kvadrat na poziciji stupca 2; svi su X.} \\
\texttt{Da, moj desni rubni stupac ima isključivo znakove X u svih pet redova.} \\
            }
        }
    }
    & & \\ \\

    \theutterance \stepcounter{utterance}  
    & & & \multicolumn{2}{p{0.3\linewidth}}{
        \cellcolor[rgb]{0.9,0.9,0.9}{
            \makecell[{{p{\linewidth}}}]{
                \texttt{\tiny{[GM$|$GM]}}
                \texttt{continue} \\
            }
        }
    }
    & & \\ \\

    \theutterance \stepcounter{utterance}  
    & & \multicolumn{4}{p{0.6\linewidth}}{
        \cellcolor[rgb]{0.9,0.9,0.9}{
            \makecell[{{p{\linewidth}}}]{
                \texttt{\tiny{[P1$\langle$GM]}}
                \texttt{Sada postavi pitanje kako bi saznao nove aspekte moje mreže koji se možda razlikuju od tvoje mreže. Započni s "PITANJE:" i ne dodavaj ništa drugo.} \\
            }
        }
    }
    & \\ \\

    \theutterance \stepcounter{utterance}  
    & \multicolumn{4}{p{0.6\linewidth}}{
        \cellcolor[rgb]{0.9,0.9,0.9}{
            \makecell[{{p{\linewidth}}}]{
                \texttt{\tiny{[P1$\rangle$GM]}}
                \texttt{PITANJE: Ima li tvoj gornji red (red 1) prazan kvadrat (▢) na poziciji stupca 4?} \\
            }
        }
    }
    & & \\ \\

    \theutterance \stepcounter{utterance}  
    & & & \multicolumn{2}{p{0.3\linewidth}}{
        \cellcolor[rgb]{0.9,0.9,0.9}{
            \makecell[{{p{\linewidth}}}]{
                \texttt{\tiny{[GM$|$GM]}}
                \texttt{continue} \\
            }
        }
    }
    & & \\ \\

    \theutterance \stepcounter{utterance}  
    & & \multicolumn{4}{p{0.6\linewidth}}{
        \cellcolor[rgb]{0.9,0.9,0.9}{
            \makecell[{{p{\linewidth}}}]{
                \texttt{\tiny{[GM$\rangle$P2]}}
                \texttt{ODGOVOR: } \\
\texttt{Ne, moj donji red nema prazan kvadrat na poziciji stupca 2; svi su X.} \\
\texttt{Da, moj desni rubni stupac ima isključivo znakove X u svih pet redova.} \\
\texttt{PITANJE: Ima li tvoj gornji red (red 1) prazan kvadrat (▢) na poziciji stupca 4?Započni svoj odgovor s "ODGOVOR:" i nemoj dodati ništa drugo.} \\
            }
        }
    }
    & \\ \\

    \theutterance \stepcounter{utterance}  
    & & & \multicolumn{4}{p{0.6\linewidth}}{
        \cellcolor[rgb]{0.9,0.9,0.9}{
            \makecell[{{p{\linewidth}}}]{
                \texttt{\tiny{[GM$\langle$P2]}}
                \texttt{ODGOVOR: Ne, moj gornji red nema prazan kvadrat na poziciji stupca 4; tamo je X.} \\
\texttt{PITANJE: Ima li tvoj donji red (red 5) prazan kvadrat (▢) na poziciji stupca 4?} \\
            }
        }
    }
     \\ \\

    \theutterance \stepcounter{utterance}  
    & & & \multicolumn{2}{p{0.3\linewidth}}{
        \cellcolor[rgb]{0.9,0.9,0.9}{
            \makecell[{{p{\linewidth}}}]{
                \texttt{\tiny{[GM$|$GM]}}
                \texttt{continue} \\
            }
        }
    }
    & & \\ \\

    \theutterance \stepcounter{utterance}  
    & & \multicolumn{4}{p{0.6\linewidth}}{
        \cellcolor[rgb]{0.9,0.9,0.9}{
            \makecell[{{p{\linewidth}}}]{
                \texttt{\tiny{[GM$\rangle$P2]}}
                \texttt{Sada postavi pitanje kako bi saznao nove aspekte moje mreže koji se možda razlikuju od tvoje mreže. Započni s "PITANJE:" i ne dodavaj ništa drugo.} \\
            }
        }
    }
    & \\ \\

    \theutterance \stepcounter{utterance}  
    & & & \multicolumn{4}{p{0.6\linewidth}}{
        \cellcolor[rgb]{0.9,0.9,0.9}{
            \makecell[{{p{\linewidth}}}]{
                \texttt{\tiny{[GM$\langle$P2]}}
                \texttt{PITANJE: Jesu li sva polja u tvojoj unutrašnjosti (redovi 2{-}4, stupci 2{-}4) prazni kvadrati (▢)?} \\
            }
        }
    }
     \\ \\

    \theutterance \stepcounter{utterance}  
    & & & \multicolumn{2}{p{0.3\linewidth}}{
        \cellcolor[rgb]{0.9,0.9,0.9}{
            \makecell[{{p{\linewidth}}}]{
                \texttt{\tiny{[GM$|$GM]}}
                \texttt{continue} \\
            }
        }
    }
    & & \\ \\

    \theutterance \stepcounter{utterance}  
    & & \multicolumn{4}{p{0.6\linewidth}}{
        \cellcolor[rgb]{0.9,0.9,0.9}{
            \makecell[{{p{\linewidth}}}]{
                \texttt{\tiny{[P1$\langle$GM]}}
                \texttt{ODGOVOR: Ne, moj gornji red nema prazan kvadrat na poziciji stupca 4; tamo je X.} \\
\texttt{PITANJE: Ima li tvoj donji red (red 5) prazan kvadrat (▢) na poziciji stupca 4?} \\
\texttt{PITANJE: Jesu li sva polja u tvojoj unutrašnjosti (redovi 2{-}4, stupci 2{-}4) prazni kvadrati (▢)?Započni svoj odgovor s "ODGOVOR:" i nemoj dodati ništa drugo.} \\
            }
        }
    }
    & \\ \\

    \theutterance \stepcounter{utterance}  
    & \multicolumn{4}{p{0.6\linewidth}}{
        \cellcolor[rgb]{0.9,0.9,0.9}{
            \makecell[{{p{\linewidth}}}]{
                \texttt{\tiny{[P1$\rangle$GM]}}
                \texttt{ODGOVOR:} \\
\texttt{Ne, moj donji red nema prazan kvadrat na poziciji stupca 4; tamo je X.} \\
\texttt{Da, sva polja u mojoj unutrašnjosti (redovi 2{-}4, stupci 2{-}4) su prazni kvadrati (▢).} \\
            }
        }
    }
    & & \\ \\

    \theutterance \stepcounter{utterance}  
    & & & \multicolumn{2}{p{0.3\linewidth}}{
        \cellcolor[rgb]{0.9,0.9,0.9}{
            \makecell[{{p{\linewidth}}}]{
                \texttt{\tiny{[GM$|$GM]}}
                \texttt{continue} \\
            }
        }
    }
    & & \\ \\

    \theutterance \stepcounter{utterance}  
    & & \multicolumn{4}{p{0.6\linewidth}}{
        \cellcolor[rgb]{0.9,0.9,0.9}{
            \makecell[{{p{\linewidth}}}]{
                \texttt{\tiny{[P1$\langle$GM]}}
                \texttt{Sada postavi pitanje kako bi saznao nove aspekte moje mreže koji se možda razlikuju od tvoje mreže. Započni s "PITANJE:" i ne dodavaj ništa drugo.} \\
            }
        }
    }
    & \\ \\

    \theutterance \stepcounter{utterance}  
    & \multicolumn{4}{p{0.6\linewidth}}{
        \cellcolor[rgb]{0.9,0.9,0.9}{
            \makecell[{{p{\linewidth}}}]{
                \texttt{\tiny{[P1$\rangle$GM]}}
                \texttt{PITANJE: Ima li tvoj lijevi rubni stupac (stupac 1) prazan kvadrat (▢) u bilo kojem retku osim u redovima 1 i 5?} \\
            }
        }
    }
    & & \\ \\

    \theutterance \stepcounter{utterance}  
    & & & \multicolumn{2}{p{0.3\linewidth}}{
        \cellcolor[rgb]{0.9,0.9,0.9}{
            \makecell[{{p{\linewidth}}}]{
                \texttt{\tiny{[GM$|$GM]}}
                \texttt{continue} \\
            }
        }
    }
    & & \\ \\

    \theutterance \stepcounter{utterance}  
    & & \multicolumn{4}{p{0.6\linewidth}}{
        \cellcolor[rgb]{0.9,0.9,0.9}{
            \makecell[{{p{\linewidth}}}]{
                \texttt{\tiny{[GM$\rangle$P2]}}
                \texttt{ODGOVOR:} \\
\texttt{Ne, moj donji red nema prazan kvadrat na poziciji stupca 4; tamo je X.} \\
\texttt{Da, sva polja u mojoj unutrašnjosti (redovi 2{-}4, stupci 2{-}4) su prazni kvadrati (▢).} \\
\texttt{PITANJE: Ima li tvoj lijevi rubni stupac (stupac 1) prazan kvadrat (▢) u bilo kojem retku osim u redovima 1 i 5?Započni svoj odgovor s "ODGOVOR:" i nemoj dodati ništa drugo.} \\
            }
        }
    }
    & \\ \\

    \theutterance \stepcounter{utterance}  
    & & & \multicolumn{4}{p{0.6\linewidth}}{
        \cellcolor[rgb]{0.9,0.9,0.9}{
            \makecell[{{p{\linewidth}}}]{
                \texttt{\tiny{[GM$\langle$P2]}}
                \texttt{ODGOVOR: Ne, moj lijevi rubni stupac nema prazan kvadrat ni u jednom retku, uključujući redove 1 i 5; svi su X.} \\
\texttt{PITANJE: Mogu li zaključiti da je tvoja mreža potpuno ispunjena X{-}ovima po cijelom vanjskom rubu (svi rubni redovi i stupci) bez ikakvih praznina, a unutrašnjost 3x3 je ispunjena isključivo praznim kvadratima?} \\
            }
        }
    }
     \\ \\

    \theutterance \stepcounter{utterance}  
    & & & \multicolumn{2}{p{0.3\linewidth}}{
        \cellcolor[rgb]{0.9,0.9,0.9}{
            \makecell[{{p{\linewidth}}}]{
                \texttt{\tiny{[GM$|$GM]}}
                \texttt{continue} \\
            }
        }
    }
    & & \\ \\

    \theutterance \stepcounter{utterance}  
    & & \multicolumn{4}{p{0.6\linewidth}}{
        \cellcolor[rgb]{0.9,0.9,0.9}{
            \makecell[{{p{\linewidth}}}]{
                \texttt{\tiny{[GM$\rangle$P2]}}
                \texttt{Sada donesi odluku. Što misliš: jesu li tvoja mreža i druga opisana mreža iste? Napiši "ODLUKA: ista mreža" ako misliš da su iste mreže ili "ODLUKA: različite mreže" ako misliš da su različite mreže. Ne dodaj ništa drugo.} \\
            }
        }
    }
    & \\ \\

    \theutterance \stepcounter{utterance}  
    & & & \multicolumn{4}{p{0.6\linewidth}}{
        \cellcolor[rgb]{0.9,0.9,0.9}{
            \makecell[{{p{\linewidth}}}]{
                \texttt{\tiny{[GM$\langle$P2]}}
                \texttt{ODLUKA: različite mreže} \\
            }
        }
    }
     \\ \\

    \theutterance \stepcounter{utterance}  
    & & & \multicolumn{2}{p{0.3\linewidth}}{
        \cellcolor[rgb]{0.9,0.9,0.9}{
            \makecell[{{p{\linewidth}}}]{
                \texttt{\tiny{[GM$|$GM]}}
                \texttt{continue} \\
            }
        }
    }
    & & \\ \\

    \theutterance \stepcounter{utterance}  
    & & & \multicolumn{2}{p{0.3\linewidth}}{
        \cellcolor[rgb]{0.9,0.9,0.9}{
            \makecell[{{p{\linewidth}}}]{
                \texttt{\tiny{[GM$|$GM]}}
                \texttt{success} \\
            }
        }
    }
    & & \\ \\

    \theutterance \stepcounter{utterance}  
    & & \multicolumn{4}{p{0.6\linewidth}}{
        \cellcolor[rgb]{0.9,0.9,0.9}{
            \makecell[{{p{\linewidth}}}]{
                \texttt{\tiny{[P1$\langle$GM]}}
                \texttt{ODGOVOR: Ne, moj lijevi rubni stupac nema prazan kvadrat ni u jednom retku, uključujući redove 1 i 5; svi su X.} \\
\texttt{PITANJE: Mogu li zaključiti da je tvoja mreža potpuno ispunjena X{-}ovima po cijelom vanjskom rubu (svi rubni redovi i stupci) bez ikakvih praznina, a unutrašnjost 3x3 je ispunjena isključivo praznim kvadratima?} \\
\texttt{Sada donesi odluku. Što misliš: jesu li tvoja mreža i druga opisana mreža iste? Napiši "ODLUKA: ista mreža" ako misliš da su iste mreže ili "ODLUKA: različite mreže" ako misliš da su različite mreže. Ne dodaj ništa drugo.} \\
            }
        }
    }
    & \\ \\

    \theutterance \stepcounter{utterance}  
    & \multicolumn{4}{p{0.6\linewidth}}{
        \cellcolor[rgb]{0.9,0.9,0.9}{
            \makecell[{{p{\linewidth}}}]{
                \texttt{\tiny{[P1$\rangle$GM]}}
                \texttt{ODLUKA: različite mreže} \\
            }
        }
    }
    & & \\ \\

    \theutterance \stepcounter{utterance}  
    & & & \multicolumn{2}{p{0.3\linewidth}}{
        \cellcolor[rgb]{0.9,0.9,0.9}{
            \makecell[{{p{\linewidth}}}]{
                \texttt{\tiny{[GM$|$GM]}}
                \texttt{continue} \\
            }
        }
    }
    & & \\ \\

    \theutterance \stepcounter{utterance}  
    & & & \multicolumn{2}{p{0.3\linewidth}}{
        \cellcolor[rgb]{0.9,0.9,0.9}{
            \makecell[{{p{\linewidth}}}]{
                \texttt{\tiny{[GM$|$GM]}}
                \texttt{success} \\
            }
        }
    }
    & & \\ \\

\end{supertabular}
}

\end{document}
